\section{Estimation}
\label{sec:estimation}
This section describes the estimation method, including data selection, priors and posteriors of parameters. The methodology follows \cite{smetsShocksFrictionsUS2007} and \cite{justinianoInvestmentShocksBusiness2010}.

\subsection{Data}
The data set contains U.S. data on GDP growth rate, inflation rate, the federal funds rate, hours worked, investment growth rate, and government spending growth rate. The data period is from 1983:Q1 to 2002:Q4. The period selection is based on three considerations. First, since the economic model assumes a constant linear growth trend, we may face the problem of changing growth rates if we use a data set over long periods. Second, the data is selected to be comparable with \cite{justinianoInvestmentShocksBusiness2010}, who used the data from 1954:Q3 to 2004:Q4. Third, the data period avoids the periods where the economy faced the zero-lower bound constraint, as the model lacks a clear picture of monetary policy in such regimes.

All data are collected from the FRED database, and are processed to represent per-capita levels because there is no detailed modeling of population growth in the baseline model. The detail of the data process is described in Appendix~\ref{app:data}.
\begin{equation}\label{eq:observation}
    \begin{bmatrix}
        \Delta \ln GDP_t \\
        \Delta \ln P_t \\
        FEDFUNDS_t\\
        \ln HOURS_t \\
        \Delta \ln INV_t \\
        \Delta \ln GOV_t \\
    \end{bmatrix}
    =
    \begin{bmatrix}
        100 \ln \gamma \\
        400\pi_{ss} \\
        400R_{ss} \\
        \bar{\ell}\\
        400\ln \gamma \\
        400\ln \gamma \\
    \end{bmatrix}
    +
    \begin{bmatrix}
        \hat{y}_t - \hat{y}_{t-1} \\
        \hat{\pi}_t \\
        \hat{R}_t \\
        \hat{\ell}_t\\
        \hat{i}_t - \hat{i}_{t-1}\\
        \hat{g}_t - \hat{g}_{t-1}
    \end{bmatrix}
    +
    \Sigma_t^{meas}
\end{equation}

Equation \eqref{eq:observation} describes the relationship between the observed data and the state variables, which is called the observation equation. $\pi_{ss}$ abd $R_{ss}$ are steady-state values of these variables, and variables with hats indicate the log-linearized variables. $\pi_{ss}$ and $R_{ss}$ are multiplied by 4 to indicate the year-to-year inflation rate and the federal funds rate. Moreover, because I use the demeaned $HOURS_t$, $\bar{\ell} = \ell_{ss} - \mu_{hours}$, where $\ell_{ss}$ is the steady-state labor input level and $\mu_{hours}$ is the mean of the data.

$\Sigma^{meas}_t$ is a $6 \times 1$ vector, containing the measurement error for each data variable. The measurement error helps the model to adjust the number of shocks flexibly, while the stochastic singularity problem may be raised with fewer shocks than the data \citep{ruge-murciaMethodsEstimateDynamic2007}. The value of the measurement error is set as 4\% of the variance of each series, meaning that there is a small part of variance in the data caused by the data collection process.

Compared with \cite{smetsShocksFrictionsUS2007} and \cite{justinianoInvestmentShocksBusiness2010}, I add the data of government spending, because there are fewer shocks in my model than theirs and the government spending shock will dominate the fluctuation otherwise. In addition, the data on consumption and wages are dropped, as more shocks and mechanisms are required for the model to fit these additional data. Nevertheless, the result of the expanded model will be examined in section~\ref{sec:robustness}.

\subsection{Prior and Posterior Distributions}
\label{sec:estimation:dist}
To facilitate the setting of the priors, I reparametrize some of the parameters. Specifically,
\begin{equation}
    \begin{aligned}
        \bar{\gamma} &= 100(\gamma - 1) \approx 100\ln\gamma,\\
        \bar{\pi} &= 400\pi_{ss},\\
        \bar{r} &= 400\left(\frac{\gamma}{\beta} - 1\right).
    \end{aligned}
\end{equation}

Accordingly, $R_ss$ is calculated as $\bar{\pi} + \bar{r}$.

Additionally, there are some parameters that suffer from weak identification problems, and therefore I set their values following \cite{smetsShocksFrictionsUS2007,leeperDynamicsFiscalFinancing2010,nireiOriginsPropagationAnimal}. The inverse Frisch elasticity in the final goods production function $\nu_p$ is set as 6, indicating that the steady-state markup is 0.12. The depreciation rate of capital $\delta$ is set as 0.025, and the steady ratio of government spending to output $s_g$ is set as $0.2$.

The details about the priors and posteriors are summarized in table \ref{tab:prior_posterior_params}, where the posterior is based on the Metropolis-Hastings algorithm with 100,000 sample size and 0.5 burn-in rate. The priors of the parameters are relatively flat so that less restrictions are imposed on the posteriors.

% Auto-generated summary table combining priors (prior.m) and posteriors (posterior_params_gov_data.tex)
% file_name: gov_data

\begin{table}[!htbp]
\centering
\caption{Prior distributions and posterior summaries (mean, std, 5\%, 95\%) for parameters in the baseline estimation.}
\label{tab:prior_posterior_params}
\begin{tabular}{c @{\hspace{2em}} l r @{\hspace{1.5em}} r c @{\hspace{1.5em}} r r r r}
\hline\hline
& \multicolumn{3}{c}{Prior} & & \multicolumn{4}{c}{Posterior} \\
\cline{2-4} \cline{6-9}
Parameter & Dist. & Mean & Std & & Mean & Std & 5\% & 95\% \\
\hline
$\sigma$ & Gamma & 2.00 & 0.50 & & 2.56 & 0.37 & 1.99 & 3.19 \\
$\sigma_\ell$ & Gamma & 2.00 & 0.50 & & 1.77 & 0.26 & 1.37 & 2.23 \\
$\kappa_p$ & Beta & 0.30 & 0.10 & & 0.43 & 0.07 & 0.32 & 0.55 \\
$\alpha$ & Normal & 0.30 & 0.05 & & 0.41 & 0.02 & 0.38 & 0.44 \\
$\theta_w$ & Beta & 0.50 & 0.20 & & 0.50 & 0.10 & 0.35 & 0.67 \\
$\bar{N}$ & Normal & 0.00 & 2.00 & & 0.96 & 1.69 & -1.82 & 3.77 \\
$\bar{\gamma}$ & Normal & 0.40 & 0.20 & & 0.33 & 0.03 & 0.27 & 0.37 \\
$\bar{r}$ & Gamma & 0.50 & 0.50 & & 1.80 & 0.44 & 1.07 & 2.49 \\
$\bar{\pi}$ & Gamma & 7.00 & 2.00 & & 4.05 & 0.69 & 3.05 & 5.36 \\
$\psi_{\pi}$ & Normal & 1.50 & 0.25 & & 1.87 & 0.18 & 1.60 & 2.18 \\
$\varphi_i$ & Gamma & 4.00 & 1.50 & & 8.35 & 1.98 & 5.40 & 11.79 \\
$\psi_u$ & Beta & 0.50 & 0.15 & & 0.09 & 0.03 & 0.04 & 0.14 \\
$\phi_g$ & Gamma & 0.07 & 0.05 & & 0.04 & 0.01 & 0.02 & 0.06 \\
$\psi_{y}$ & Gamma & 0.50 & 0.25 & & 0.28 & 0.13 & 0.10 & 0.52 \\
$\rho_a$ & Beta & 0.60 & 0.20 & & 0.94 & 0.02 & 0.89 & 0.98 \\
$\rho_g$ & Beta & 0.60 & 0.20 & & 0.99 & 0.00 & 0.99 & 0.99 \\
$\rho_R$ & Beta & 0.60 & 0.20 & & 0.76 & 0.04 & 0.70 & 0.82 \\
$\rho_s$ & Beta & 0.60 & 0.20 & & 0.99 & 0.01 & 0.97 & 1.00 \\
$\rho_d$ & Beta & 0.60 & 0.20 & & 0.87 & 0.04 & 0.80 & 0.93 \\
$\sigma_a$ & InvGamma & 0.50 & 4.00 & & 0.45 & 0.06 & 0.36 & 0.56 \\
$\sigma_g$ & InvGamma & 0.50 & 4.00 & & 0.89 & 0.07 & 0.78 & 1.01 \\
$\sigma_R$ & InvGamma & 0.50 & 4.00 & & 0.24 & 0.03 & 0.20 & 0.29 \\
$\sigma_s$ & InvGamma & 0.50 & 4.00 & & 1.41 & 0.35 & 0.89 & 2.04 \\
$\sigma_d$ & InvGamma & 0.50 & 4.00 & & 1.35 & 0.16 & 1.10 & 1.62 \\
\hline
\end{tabular}
\end{table}


$\varphi_i$ is the parameter in investment adjustment function, which is defined in equation~\eqref{eq:loglin_inv} in Appendix~\ref{app:model}. $\psi_u$ is the parameter in the capital utilization adjustment cost function defined in equation~\eqref{eq:loglin_u}. Moreover, $\kappa_p$ is the coefficient in the Phillips curve, defined in equation~\eqref{eq:nkpc}.

There are a few posteriors worth to be examined. The first one is $\psi_u$, whose posteriors are quite different among literature. \cite{smetsShocksFrictionsUS2007} estimates $\psi_u$ about $0.5$,\cite{justinianoInvestmentShocksBusiness2010} gives about $0.8$, and Section~\ref{sec:robustness} gives about $0.1$, even though the model structure is quite similar. The second one is $\rho_s$, as it is pretty much to the unit root. This extreme value may be a result of lacking enough shocks to interpret the volatility in the data, as the result Section~\ref{sec:robustness} shows a more temperate persistence in IST shock. Finally, $kappa_p$ and $\theta_w$ are around $0.5$, which indicates the Calvo parameter is near $0.5$ for the final goods, and $0.2$ for the wage. This means $50\%$ of final goods producers can reset their price, and $80\%$ of households can reset their wage. I flag these problematic posteriors here and will investigate their implications for other results in the next section.